\section{Persistence and packing}
\label{sec:persistence}

We now bound the average of $N$ at the terminal level. The stopped
descent estimate controls the contribution of one excursion of high
growth, but a path may enter several excursions. We will bound the
total area of their entry squares. The argument has two parts: an
entry produces a nearby region on which few later entries can occur,
and a geometric packing estimate converts these regions into the
required area bound.

Fix the root square $Q_0$ and the constants $A,C_0,C_1$ from
Section~\ref{sec:subdivision}. We consider finite subdivision trees whose
terminal squares have a common side length and for which $K_Q\ge a_0>0$
at every level. All constants in this section may depend on these fixed
data, including $a_0$, but are independent of the eigenfunction and the
depth of the tree. Let $L$ be the threshold in
Lemma~\ref{lem:descent}, and let $P,\sigma$ be the constants in
Proposition~\ref{prop:stability}. Proposition~\ref{prop:initial-growth}
and the fixed coordinate volume comparisons give a constant
$C_{\mathrm{init}}$ such that $N(Q_0)\le C_{\mathrm{init}}$.

\subsection{Entries and persistent nearby regions}

\begin{definition}[Excursions]
\label{def:excursions}
Choose a threshold $H>\max\{L,C_{\mathrm{init}}\}$. Along each descending
grid path, an excursion starts the first time $N(Q)>H$ and ends the first
subsequent time $N(Q)\le L$. After an excursion has ended, apply the same
rule again. A square at which an excursion starts is called an
\emph{entry}. An excursion is \emph{active} at a square if it has started
and has not yet ended there.
\end{definition}

These decisions depend only on a square and its ancestors, so each entry
is a whole grid square. Write $\mathcal E$ for the family of entries.
The root is not an entry. In addition to the
lower bound defining an entry, we have a uniform upper bound on its
growth.

\begin{lemma}[Growth at an entry]
\label{lem:entry-bound}
For all sufficiently large fixed $H$, every entry $Q$ and its grid
parent $Q^p$ satisfy
\begin{equation}
\label{eq:entry-bound}
 N(Q^p)\le H,
 \qquad
 H<N(Q)\le
 A^2\bigl[\tfrac12+(R+1)C_1\bigr]
       \left(H+\frac{C_0}{a_0}\right).
\end{equation}
\end{lemma}

\begin{proof}
Immediately before an entry the path is inactive. An inactive square
cannot have $N>H$, since that would start an excursion, so
$N(Q^p)\le H$. Apply Proposition~\ref{prop:subdivision} at $Q^p$ with
$b=0$ and $S=HK_{Q^p}+C_0$. Indeed, $E(Q^p,0)\le HK_{Q^p}\le S$, and
\[
 \frac{K_{Q^p}}{S}\le\frac1H,
 \qquad
 \frac{K_{Q^p}^2s(Q^p)}{S^2}\le\frac1{H^2},
\]
since $s(Q^p)\le1$. Thus
\eqref{eq:subdivision-hypotheses} holds when $H$ is sufficiently large.
All terms in the averages in \eqref{eq:subdivision-estimate} are
nonnegative. There are $A^2$ children, so the tilt $b_Q$ assigned to the
particular child $Q$ satisfies
\[
 m(Q,b_Q)\le \frac{A^2S}{2K_{Q^p}},
 \qquad
 |b_Q|\le \frac{A^2C_1S}{K_{Q^p}}.
\]
Combining these inequalities with
\eqref{eq:untilted-comparison} and $K_{Q^p}\ge a_0$ proves the upper
bound in \eqref{eq:entry-bound}.
\end{proof}

For a square $Q$ and a number $D\ge1$, write $DQ$ for its concentric
dilation by the factor $D$. The following proposition supplies a region
of controlled future entry count near each entry. Its location need not
be inside the entry square.

\begin{proposition}[Persistence near an entry]
\label{prop:persistence}
There are a finite threshold $H$, a number $\delta>0$, and an integer
$J\ge0$ such that every entry $Q$ admits a measurable set
$Y_Q\subset D_*Q$, where $D_*=20A$, with
\begin{equation}
\label{eq:persistence}
 \begin{split}
 |Y_Q|&\ge\delta|Q|,\\
 \#\{F\in\mathcal E:\ s(F)<s(Q),\ y\in F\}&\le J
 \qquad\text{for almost every }y\in Y_Q.
 \end{split}
\end{equation}
Only entries of the given tree are counted, including when $Y_Q$
extends outside the root.
\end{proposition}

\begin{proof}
The contradiction argument below extracts a nonconstant profile,
finds a nearby square with small affine excess relative to a nonzero
slope, and transfers that square to Proposition~\ref{prop:stability}.

We may require $H$ to exceed any fixed threshold needed above. Suppose
there is no triple $(H,\delta,J)$ with the asserted uniform property.
Choose $H_j\to\infty$, all above these fixed thresholds and at least
$2$, and apply this negation with
$\delta_j=1/H_j$ and $J_j=\lceil H_j\rceil$. For each $j$ there are an
eigenfunction, a finite tree, and an entry $Q_j$ for which
\eqref{eq:persistence} fails with these parameters.

There must be more than $J_j$ levels strictly below $Q_j$. Otherwise the
finer-entry count would be at most $J_j$ everywhere, and $Y_{Q_j}=Q_j$
would satisfy \eqref{eq:persistence}. In particular, every fixed number
of further levels is eventually available.

Rescale about the center $x_j$ of the parent $Q_j^p$, using its side
$s_j=s(Q_j^p)$ as unit length. Write $k_j$ for the eigenfrequency of
$u_j$, and put
\[
 K_j=k_js_j,\qquad S_j=H_jK_j,\qquad
 U_j(\xi)=u_j(x_j+s_j\xi),\qquad
 p_j(\xi)=p(x_j+s_j\xi).
\]
The rescaled equation is
$\Delta U_j+K_j^2p_jU_j=0$. Its coefficients satisfy
$\|p_j\|_\infty\le C$ and $\|\nabla p_j\|_\infty\le Cs_j$.
Moreover,
\[
 S_j\ge H_ja_0\longrightarrow\infty,
 \qquad \frac{K_j}{S_j}=\frac1{H_j}\longrightarrow0,
 \qquad
 \frac{K_j^2\|\nabla p_j\|_\infty}{S_j^2}
 \le\frac{Cs_j}{H_j^2}\longrightarrow0.
\]
Normalize $|U_j|^2\,d\xi$ to have mass one on $\D_R$.
Lemma~\ref{lem:profile-extraction} and
Theorem~\ref{thm:subharmonic-profile}, with zero tilt, give a
subharmonic profile $V$ on $\D_R$ satisfying
\[
 V\le0,\qquad \sup_{\overline{\D_1}}V\ge-1.
\]
For the second assertion, the parent inequality gives
$E(Q_j^p,0)/S_j\le1$, and the compact upper bound in
\eqref{eq:profile-bounds} applies to the closed unit disk. Thus $V$ is
not identically $-\infty$. Its supremum on every nonempty open disk in
$\D_R$ is therefore finite.

There are only $A^2$ possible positions of the entry inside its parent.
After a subsequence its rescaled center is a fixed point $y_e$. The
entry inequality and the compact upper and open lower bounds in
\eqref{eq:profile-bounds} imply
\[
 \frac1A
 \le \liminf_j\frac{E(Q_j,0)}{S_j}
 \le \limsup_j\frac{E(Q_j,0)}{S_j}
 \le
 \sup_{\overline{\D_{R/A}(y_e)}}V
       -\sup_{\D_{1/A}(y_e)}V.
\]
The finiteness of the second supremum justifies the subtraction of
these bounds. Since $A>10R$, both disks are contained in $\D_2$.
Consequently $V$ is nonconstant on $\D_2$.

The representation in Lemma~\ref{lem:potential-decomposition} implies
$V\in W^{1,1}(\D_2)$: the gradient of the logarithmic kernel is locally
integrable in two dimensions, and the representing measure has finite
mass. There is a Lebesgue point $z\in\D_2$ of $\nabla V$ for which
$c=\nabla V(z)\ne0$. Indeed, otherwise the weak gradient would vanish
almost everywhere on the connected disk, forcing $V$ to be constant
almost everywhere; its subharmonic representative, which is determined
by local averages, would then be constant everywhere.

Extend the fine subdivision grids in the rescaled parent coordinates
to the whole plane. For $r=A^{-m}$, choose a grid square with center
$y_m$ such that $|y_m-z|\le r$. These grids have edges
$-\tfrac12+nA^{-m}$ in each coordinate and are the same for every $j$.
We claim that
\begin{equation}
\label{eq:affine-lebesgue}
 \sup_{\overline{\D_{Rr}(y_m)}}(V(x)-c\cdot x)
 -\sup_{\D_r(y_m)}(V(x)-c\cdot x)=o(r)
 \qquad(m\longrightarrow\infty).
\end{equation}
To prove this, let $W(x)=V(x)-c\cdot x$, and let $a_r$ be its average
on $B_r=\D_{(2R+2)r}(z)$. For large $m$ these disks are contained in
$\D_2$. Poincare's inequality and the Lebesgue-point property give
\[
 \av_{B_r}|W-a_r|
 \le C_Rr\,\av_{B_r}|\nabla V-c|=o(r).
\]
If $x\in\overline{\D_{Rr}(y_m)}$, the disk
$\D_{(R+1)r}(x)$ lies in $B_r$. The submean inequality for the
subharmonic function $W$ therefore gives, uniformly in such $x$,
\[
 W(x)\le a_r+4\av_{B_r}|W-a_r|=a_r+o(r).
\]
On the other hand,
\[
 \sup_{\D_r(y_m)}W
 \ge\av_{\D_r(y_m)}W
 \ge a_r-(2R+2)^2\av_{B_r}|W-a_r|=a_r-o(r).
\]
This proves \eqref{eq:affine-lebesgue}. Subharmonicity is essential
here: it turns the average control into a uniform upper bound.

Now fix $m\ge2$ sufficiently large that $Rr\le1/4$ and the quotient of
the left side of \eqref{eq:affine-lebesgue} by $r|c|$ is strictly less
than $\sigma$. In the original coordinates let $I_j$ be the square
with center $x_j+s_jy_m$ and side $s_jr$, and assign it the tilt
$b_{I_j}=H_jc$. The additional weight in parent coordinates is
$e^{-2S_jc\cdot\xi}$, up to a constant canceled in the norm ratio.
The weighted bounds of Lemma~\ref{lem:profile-extraction} therefore yield
\begin{equation}
\label{eq:persistence-transfer}
 \limsup_j\frac{m(I_j,b_{I_j})}{|H_jc|}
 \le
 \frac{
 \sup_{\overline{\D_{Rr}(y_m)}}(V(x)-c\cdot x)
 -\sup_{\D_r(y_m)}(V(x)-c\cdot x)}{r|c|}
 <\sigma.
\end{equation}
The factor is exactly $1/(r|c|)$, since $K_{I_j}=rK_j$ and
$S_j=H_jK_j$. Because $c\ne0$, we also have $|H_jc|\ge P$ for all
large $j$. Thus \eqref{eq:stability-hypotheses} holds at $I_j$.

We verify that Proposition~\ref{prop:stability} can be used on the
whole subtree of $I_j$ down to the original bottom. The chosen $m$ is
fixed before $j$ tends to infinity, so that level is eventually
available. All squares of a given extended-grid level have the same
physical scale as the original tree, and hence have $K\ge a_0$.
The center $y_m$ stays in $\D_3$, and $Rr\le1/4$. The disks required
for $I_j$ and all its descendants are consequently contained in
$\D_{5s_j}(x_j)$: even the root margin has rescaled radius $2Rr$,
and $|y_m|+2Rr<3+1/2<5$. Since $Q_j^p\subset Q_0$, we have
$|x_j-x_{Q_0}|\le s_0/\sqrt2$ and $s_j\le s_0$, so this disk lies in
$\D_{6s_0}(x_{Q_0})\subset\D_{2Rs_0}(x_{Q_0})$, within the fixed
coordinate patch. This also covers the larger disk margins required
when regarding $I_j$ as the root of a subtree.

Proposition~\ref{prop:stability} now gives a measurable subset of
$I_j$ of area at least $|I_j|/2$ on which every visited square from
$I_j$ through the bottom has $N>L$. The extended grids align with the
original root: a root edge differs from the corresponding parent edge
by an integer multiple of $s_j$, and the extended parent grid has
rescaled edges $-1/2+nA^{-m}$. Thus a fine grid square is either
contained in $Q_0$ or has interior disjoint from it. In the latter case
its points belong to no
entry of the original tree. In the former case its descendants are
exactly those of the original tree. From the level of $I_j$ onward
there can be at most one entry along any of these paths, since two
starts require an intervening value $N\le L$. Before that level there
are at most $m$ relevant levels below the tested entry. Thus the number
of strictly finer entries is at most $m+1$ on this subset.

Finally, the subset lies in $20A Q_j$: in parent units that dilated
square has half-side $10$, whereas $y_m\in\D_3$ and the entry center
lies in the unit parent square. Its area relative to the entry is at
least
\[
 \frac{|I_j|}{2|Q_j|}=\frac12 A^{2-2m}>0.
\]
Both this positive number and $m+1$ are fixed independently of $j$.
They eventually satisfy the tested requirements
$\delta_j=1/H_j$ and $J_j=\lceil H_j\rceil$, contradicting the choice
of $Q_j$. This proves the proposition.
\end{proof}

\subsection{Packing with witnesses outside the squares}

The next geometric lemma converts the nearby regions just constructed
into a bound on the sum of entry areas. We state it for an arbitrary
finite family in a subdivision tree.

\begin{lemma}[Packing from nearby witnesses]
\label{lem:packing}
Let $\mathcal F$ be a finite family of squares in the $A$-adic
subdivision tree of a square $Q_0$, where $A\ge2$. Suppose there are
fixed $D\ge1$, $\delta>0$, and an integer $J\ge0$ such that every
$Q\in\mathcal F$ admits a measurable set $Y_Q\subset DQ$ satisfying
\[
 |Y_Q|\ge\delta|Q|,
 \qquad
 \#\{F\in\mathcal F:s(F)<s(Q),\ y\in F\}\le J
 \quad\text{for almost every }y\in Y_Q.
\]
Then
\begin{equation}
\label{eq:packing}
 \sum_{Q\in\mathcal F}|Q|
 \le C(A,D,\delta,J)|Q_0|.
\end{equation}
\end{lemma}

\begin{proof}
Ignore grid boundaries and the boundaries of the finitely many dilates
$DQ$ with $Q\in\mathcal F$; these form a null set. Define
\[
 n(Q)=1+\#\{F\in\mathcal F:Q\subsetneq F\}.
\]
Squares of a fixed rank $n$ have disjoint interiors. At each rank,
process the squares in decreasing order of side length, retaining a
square if its $D$-dilate has interior disjoint from the previously
retained dilates. Let $\mathcal S$ be the union of these selected
families. Every discarded square $Q$ has its dilate meeting that of a
selected square $P$ with $s(P)\ge s(Q)$; hence
$Q\subset(2D+1)P$. Covering the original disjoint squares at that rank
by these larger dilates and then summing over ranks gives
\[
 \sum_{Q\in\mathcal S}|Y_Q|
 \ge\delta\sum_{Q\in\mathcal S}|Q|
 \ge\alpha\sum_{Q\in\mathcal F}|Q|,
 \qquad \alpha=\frac{\delta}{(2D+1)^2}.
\]
The selected witness sets at a single rank are disjoint up to null
sets. It remains to control overlap between different ranks. For a
point $y$ outside the ignored boundaries, put
\[
 n_{\mathrm{tot}}(y)=\#\{F\in\mathcal F:y\in F\}.
\]
If a witness point $y$ lies inside its square $Q$, every larger square
containing $y$ is an ancestor of $Q$, so the witness assumption makes
$n_{\mathrm{tot}}(y)$ differ from $n(Q)$ by at most $J$. An exterior
witness point can also disagree with the ancestry of $Q$ at larger
scales. For $Q\in\mathcal S$ and $y\in Y_Q$, let
\[
 d_Q(y)=
 \sum_{\substack{F\in\mathcal F\\s(F)>s(Q)}}
 \left|\mathbf 1_F(y)-\mathbf 1_{\{Q\subset F\}}\right|
\]
count the mismatches between membership of $y$ in a larger square and
that square's being an ancestor of $Q$. If $b$ is the number of larger
squares in $\mathcal F$ containing $y$, then
\[
 |n(Q)-1-b|\le d_Q(y).
\]
At most one square of side $s(Q)$ contains $y$, and the witness
assumption bounds the number of smaller squares containing it by $J$.
Decomposing $n_{\mathrm{tot}}(y)$ into these three size ranges gives
\begin{equation}\label{eq:rank-mismatch}
 |n(Q)-n_{\mathrm{tot}}(y)|\le d_Q(y)+J+1
 \quad\text{for almost every }y\in Y_Q.
\end{equation}
Thus a uniform bound for $d_Q(y)$ restricts the ranks that can contribute
at $y$ to a fixed interval. We next show that the total witness area
weighted by $d_Q$ is controlled:
\begin{equation}
\label{eq:mismatch}
 \sum_{Q\in\mathcal S}\int_{Y_Q}d_Q(y)\,dy
 \le B\sum_{F\in\mathcal F}|F|,
\end{equation}
where $B$ depends only on $A,D$.

Fix $F\in\mathcal F$ and a finer side length
$t=A^{-\ell}s(F)$, $\ell\ge1$. The grid is laminar: each square $Q$ of
side $t$ is either contained in $F$ or interior-disjoint from it. If it
can contribute a mismatch, a point of $DQ$ lies on the opposite side
of $\partial F$ from $Q$. It follows that the whole square $Q$ lies
in the collar
\[
 \{x:\dist_\infty(x,\partial F)\le(D+1)t\},
\]
where $\dist_\infty$ is distance for the maximum norm on $\R^2$.
Figure~\ref{fig:boundary-mismatch} shows the two possible directions
of this mismatch near one face of $F$.
This collar has area at most
\[
 8(D+1)s(F)t+4(D+1)^2t^2.
\]
All squares of side $t$, across all ranks, have disjoint interiors.
Since $|Y_Q|\le D^2|Q|$, their total contribution for this fixed $F$
and $t$ is at most $D^2$ times that collar area. Summing the geometric
series over $\ell\ge1$ and then summing over $F$ proves
\eqref{eq:mismatch}. For example, one may take
\[
 B=D^2\left(\frac{8(D+1)}{A-1}
            +\frac{4(D+1)^2}{A^2-1}\right).
\]

\begin{figure}[tbp]
\centering
\setlength{\unitlength}{0.9mm}
\begin{picture}(156,57)
\small
\put(35,54){\makebox(0,0){$Q\subset F,\quad y\notin F$}}
\put(112,54){\makebox(0,0){$Q\cap F=\varnothing,\quad y\in F$}}
\color{gray!35}
\put(3,13){\line(0,1){31}}
\put(67,13){\line(0,1){31}}
\put(80,13){\line(0,1){31}}
\put(144,13){\line(0,1){31}}
\color{black}
\put(35,13){\line(0,1){31}}
\put(112,13){\line(0,1){31}}
\put(35,47){\makebox(0,0){$\partial F$}}
\put(112,47){\makebox(0,0){$\partial F$}}
\put(11,42){\makebox(0,0){$F$}}
\put(88,42){\makebox(0,0){$F$}}
\put(16,13){\dashbox{1}(24,24){}}
\put(108,13){\dashbox{1}(24,24){}}
\put(20,40){\makebox(0,0){$DQ$}}
\put(128,40){\makebox(0,0){$DQ$}}
\put(24,21){\framebox(8,8){$Q$}}
\put(116,21){\framebox(8,8){$Q$}}
\put(38,25){\circle*{1.3}}
\put(40,25){\makebox(0,0)[l]{$y$}}
\put(110,25){\circle*{1.3}}
\put(108,25){\makebox(0,0)[r]{$y$}}
\put(24,18){\vector(1,0){8}}
\put(32,18){\vector(-1,0){8}}
\put(28,15){\makebox(0,0){$t$}}
\put(116,18){\vector(1,0){8}}
\put(124,18){\vector(-1,0){8}}
\put(120,15){\makebox(0,0){$t$}}
\put(35,9){\vector(1,0){32}}
\put(67,9){\vector(-1,0){32}}
\put(51,5){\makebox(0,0){$(D+1)t$}}
\put(112,9){\vector(1,0){32}}
\put(144,9){\vector(-1,0){32}}
\put(128,5){\makebox(0,0){$(D+1)t$}}
\end{picture}
\caption{Ancestry and spatial membership can disagree in either
direction. Each panel shows one witness point $y\in Y_Q\subset DQ$
and a fine square of side $t=s(Q)$ near a face of the larger square
$F$. The dashed square is $DQ$; the pale lines mark the boundary of
the collar of width $(D+1)t$ on each side of the face.}
\label{fig:boundary-mismatch}
\end{figure}

Choose an integer $M\ge1$ with $M\ge2B/\alpha$, and set
$Z_Q=\{y\in Y_Q:d_Q(y)\le M\}$. Remove also the null sets where the
witness assumption fails. Markov's inequality and
\eqref{eq:mismatch} give
\[
 \sum_{Q\in\mathcal S}|Z_Q|
 \ge\frac\alpha2\sum_{Q\in\mathcal F}|Q|.
\]
For $y\in Z_Q$, Equation~\eqref{eq:rank-mismatch} now gives
\[
 |n(Q)-n_{\mathrm{tot}}(y)|\le M+J+1.
\]
Only $2(M+J+1)+1$ ranks can therefore contribute at $y$. At a given
rank the selected witness sets are disjoint, so this is also a bound
on the total overlap of the $Z_Q$.

For $D\ge1$, the inclusions $Q\subset Q_0$ imply $DQ\subset DQ_0$.
Integrating the overlap bound over this common region yields
\[
 \frac\alpha2\sum_{Q\in\mathcal F}|Q|
 \le\bigl(2(M+J+1)+1\bigr)D^2|Q_0|.
\]
This proves \eqref{eq:packing}.
\end{proof}

Apply Lemma~\ref{lem:packing} to the family of entries, using
Proposition~\ref{prop:persistence} with $D=D_*$. We obtain a constant
$C$ such that the sum of all entry areas is at most $C|Q_0|$.

\subsection{The terminal mean growth}

\begin{theorem}[Bounded mean growth]
\label{thm:mean-growth}
For the finite trees considered in this section, let
$\mathcal T_{\mathrm{term}}$ be the collection of terminal squares and
let $s_{\mathrm{term}}$ be their common side length. There is a
constant $C$, independent of the eigenfunction and the number of
levels, such that
\begin{equation}
\label{eq:mean-growth}
 \av_{q\in\mathcal T_{\mathrm{term}}}N(q)\le C,
 \qquad
 \av_{q\in\mathcal T_{\mathrm{term}}}E(q,0)
      \le Ck s_{\mathrm{term}}.
\end{equation}
\end{theorem}

\begin{proof}
Fix the threshold $H$ supplied by
Proposition~\ref{prop:persistence}. A terminal path on which no
excursion is active has terminal growth $N\le H$. Every other terminal
path has exactly one final active entry. For an entry $Q$, let
$\Omega_Q\subset Q$ be the union of terminal squares on which that
entry's excursion is still active. These sets are disjoint as $Q$
varies.

Starting separately at each entry $Q$, set $b_Q=0$ and perform the
descent of Lemma~\ref{lem:descent}, stopping at the first square with
$N\le L$ or at the bottom, whichever comes first. Every continuing
square has $N>L$, so the descent is available regardless of its current
tilt. The low-growth stopping square, if any, is included as the last
produced value; no further step is needed there. These separate choices
of tilts do not affect excursion membership, which is defined entirely
by the untilted quantity $N$.

Along a random descendant path from $Q$, denote the stopping time by
$T$ and the successive values by $m_i,b_i$. Since $b_0=0$,
\eqref{eq:untilted-comparison} gives
\[
 N_T\le m_T+(R+1)|b_T|
 \le m_T+(R+1)\sum_{i<T}|b_{i+1}-b_i|.
\]
The uniform distribution on paths agrees with normalized planar area
on $Q$. Lemma~\ref{lem:stopped-descent} and
Lemma~\ref{lem:entry-bound} therefore imply
\[
 \begin{split}
 \int_{\Omega_Q}N_{\mathrm{term}}(x)\,dx
 &\le |Q|\,\E\left[m_T+(R+1)
                    \sum_{i<T}|b_{i+1}-b_i|\right]\\
 &\le C|Q|\bigl(N(Q)+1\bigr)
 \le C'|Q|.
 \end{split}
\]
Here $N_{\mathrm{term}}$ is constant on each terminal square, with
value $N$ of that square. The first inequality merely restricts a
nonnegative integral to the paths still active at the bottom; it does
not condition on that event or divide by its probability.

Sum over all entries and use \eqref{eq:packing}, then include the
inactive terminal paths. This gives
\[
 \int_{Q_0}N_{\mathrm{term}}(x)\,dx
 \le H|Q_0|+C'\sum_{Q\text{ entry}}|Q|
 \le C|Q_0|.
\]
Equal terminal areas turn this integral bound into the first inequality
of \eqref{eq:mean-growth}. The second follows because every terminal
square has $K_q=ks_{\mathrm{term}}$ and $E(q,0)=K_qN(q)$.
\end{proof}
